\section{Introduction}
The Ekman spiral is a famous consequence of the Coriolis force in meteorology and oceanography. Due to the interaction between shear forces in the atmosphere and ocean, the wind high up in the atmosphere, and wind over the ocean, will not give rise to a flow in the same direction at the surface. The adjustment of the flow from the surface will thus spiral, and is therefore properly named the Ekman spiral after its discoverer \citet{ekman1905}. The surface flow is often at an angle from the flow aloft by an angle of more or less \ang{45} \citep{ekman1905}. Due to this turning of the wind, the net flux of the wind will be perpendicular to the wind aloft. This is known as the Ekman transport and is crucial for studying both meteorology and oceanography.

However, the wind turning \ang{45} assumes a constant viscosity in the medium. This is unrealistic when a turbulent boundary layer is present. The boundary layer is around \SI{1}{}--\SI{2}{\kilo\meter} thick in the atmosphere and is represented by the presence of turbulence \citep{wallace2006}. This turbulence arises primarily from solar heating, causing instability and turnover, but also from mechanical turbulence from the mean wind decreasing closer to the ground due to shear force caused by frictional drag. Although turbulence is non-linear, statistical analysis can provide useful first-order approximations, such as the downward flux being negatively and linearly proportional to the local gradient. That is, one cd the sheaan define the eddy viscosity, $\nu$, for the boundary layer.

According to the mixing length hypothesis, this eddy viscosity should grow with the distance from the surface as the characteristic scales of the eddies grow. This can furthermore be described by a dimensional analysis of the boundary layer, where a logarithmic relation between the wind speed and height can be obtained \citep{wallace2006}. This showcases that the eddies, and thus the eddy viscosity, grow with height in the boundary layer. However, after a certain altitude, the turbulent boundary layer ceases to exist, and laminar flow takes over, which decreases the eddy viscosity. One can thus assume that the eddy viscosity should increase initially from the surface and then die down at a higher altitude.

Thus, in reality the viscosity should vary with altitude. This would change the wind-turning from \ang{45}, and thus also the Ekman transport. In this article, a varying viscosity with altitude will be studied under the Ekman equation. This will be done both analytically and numerically and will also be compared to model data. This will give a deeper insight into the Ekman spiral and provide insight into when a surface angle different from \ang{45} can be observed.